%% RESUMO NA LINGUA VERNACULA (elemento obrigatorio -- NBR 6028)
%%
%% Regras:
%%   - paragrafo unico, sem recuo, justificado;
%%   - Times New Roman 12, espacamento simples;
%%   - entre 150 e 500 palavras;
%%   - verbo na voz ativa e na terceira pessoa do singular;
%%   - as palavras-chave sao definidas em config/dados.tex.
%%
%% Versao publica (vitrine), sem pergunta, hipoteses e resultados.

\begin{resumo}
A pesquisa investiga a aplicação da inteligência artificial à triagem de
intimações judiciais repetitivas na Advocacia Pública, tomando como estudo de
caso o Núcleo de Cumprimento de Sentença (NCS) da Procuradoria-Geral do
Distrito Federal (PGDF), unidade especializada em cumprimentos individuais
decorrentes de julgados coletivos. Parte-se da constatação de que, em demandas
repetitivas e de baixa complexidade decisória, a definição da providência
cabível tende a ser rápida para o profissional experiente, concentrando-se o
custo operacional no tratamento individual de cada ocorrência. Nesse fluxo, a
inteligência artificial pode auxiliar na classificação da intimação e na
sugestão das providências cabíveis, enquanto mecanismos de automação podem
assumir as etapas operacionais delas decorrentes. A fundamentação teórica
articula a literatura sobre governo digital e inovação no setor público, a
transição da burocracia de balcão para a burocracia de sistema e a governança
da inteligência artificial na Administração Pública, com atenção à supervisão
humana, à transparência e à proteção de dados pessoais. A pesquisa adota
natureza aplicada e abordagem mista, instrumentalizada por estudo de caso
único selecionado por critério de tipicidade. A conformidade da solução é
examinada à luz do marco normativo aplicável à PGDF, composto pela Lei Geral
de Proteção de Dados, pela Portaria PGDF nº 160/2025 e pelo Decreto Distrital
nº 48.901/2026. Pretende-se oferecer um framework metodológico potencialmente
replicável por outras procuradorias, preservada a responsabilidade humana pelo
ato decisório e pela prática dos atos processuais.
\end{resumo}
